\chapter{Introduction}\label{ch:introduction}
\section{Motivation and application context}
Stitch begins with an application idea: a social feed whose organising surface is the physical world. A visitor could leave a photograph on a building facade, attach a note to a meeting point, or place a virtual object beside a familiar landmark. A subsequent visitor would recover that content by pointing a camera at the same scene. The repository describes this ambition through the metaphor of a social wall distributed across real surroundings \citep{stitchrepo}. This thesis uses ``Facebook in the real world'' as a motivating description, rather than as a statement of affiliation or of implemented social-network functionality.

The idea creates a correspondence problem before it creates a graphics problem. The system must decide which stored place, if any, is visible. It must then determine how content associated with that place should project into the current camera image. A visually similar building is insufficient evidence for rendering a post on a particular window. Likewise, retrieving the correct street does not identify a stable coordinate frame for a three-dimensional object. Recognition and geometric alignment therefore require separate outputs and separate evaluation criteria.

The original software provides a compact demonstration of this distinction. It detects SIFT features in a reference image and a camera frame, matches their descriptors, estimates a homography, and warps an overlay into the frame. The latest script adds optical-flow tracking and smoothing. These operations are useful for a locally planar target already represented by a reference photograph. They do not by themselves provide a place database, a city-scale search index, or a persistent world coordinate system. The engineering task is to extend that working demonstration while retaining its inspectable geometric basis.

An appearance change challenges the first correspondence step. Daylight produces different shadows, contrast and visible detail from artificial illumination at night. Seasonal vegetation may remove or reveal structures. Weather can introduce reflections or reduce contrast. These changes interact with viewpoint: a correctly recognised scene may have little shared visible surface with the reference. Long-term visual-localisation benchmarks explicitly examine these difficulties \citep{sattler2018}. A method suitable for Stitch must therefore be assessed on more than an unchanged photograph or a visually convincing overlay video.

\section{Problem statement}
Let a place database contain reference images $\mathcal{D}=\{I_i\}$, place identifiers $p_i$, and content records $\mathcal{C}_{p_i}$. Given a query image $I_q$, the system should return a place hypothesis, evidence supporting that hypothesis, a geometric transformation where identifiable, and the associated content. A rejection is a valid result when the evidence is inadequate. The planar research prototype studied here returns
\begin{equation}
 \mathcal{F}(I_q,\mathcal D)=\bigl(\widehat p,\widehat H_{r\rightarrow q},\mathcal C_{\widehat p},s\bigr),
\end{equation}
where $s$ contains diagnostic matching and verification information. When no candidate passes verification, $\widehat p$ and $\widehat H$ are absent and no anchored content is rendered.

The desired invariance is conditional. A representation should ignore exposure and other nuisance changes while retaining evidence that distinguishes places. It cannot discard all local structure and still determine the intended surface. Nor can it restore image information destroyed by underexposure. ``Appearance-invariant'' in the project title therefore identifies the research objective; it does not claim mathematical invariance or universal robustness. The experiments measure how particular pretrained representations behave under documented changes.

The resource constraint is also explicit. The available host is a laptop with an Apple M2 processor and a limited-memory configuration, recorded in the environment artifact. The implementation uses CPU inference because the recorded runtime did not expose an available accelerator. A mobile-class design is investigated through compact descriptors, bounded feature counts, model export and quantisation. A laptop CPU timing cannot establish performance on a phone, and the thesis does not make that substitution.

\section{Objectives}
The first objective is a reproducible classical reference. The original scripts must be inspected and executed on their bundled video without requiring an unavailable private camera stream. A refactored implementation should retain the SIFT--FLANN--RANSAC structure while making image coordinates, failures and rendering masks explicit. This baseline anchors the discussion in measured behaviour rather than an assumed weakness of handcrafted features.

The second objective is a controlled comparison of learned components. SuperPoint with LightGlue supplies local correspondences; CosPlace supplies an image-level descriptor. Each has a different role. Their combination is evaluated alongside exhaustive local matching and global-only retrieval. Identical query lists and geometry settings are used within each protocol so that a benefit is not explained by an easier test subset or a more permissive RANSAC threshold.

The third objective is to connect recognition to content. A place record must support reference photographs and posted assets, and a query must retrieve the content associated with its verified place. The implemented output is an image overlay attached to a reference-plane polygon. Video capture and rendering are included in the demonstration. The more ambitious multi-user service and three-dimensional anchor lifecycle are specified as architectural extensions, with implementation boundaries stated explicitly.

The fourth objective is a thesis whose evidence can be audited. Raw predictions and pair statistics are retained. Plotting and table generation consume those files, and prose-level headline values use generated LaTeX macros. Dataset and checkpoint hashes identify the inputs. Commands and exit codes identify the actions that produced results, including failed download and environment checks. The report should make an unmeasured quantity visibly unavailable rather than filling a requested table with an invented estimate.

\section{Research questions}
\begin{enumerate}
\item Does pretrained local matching improve correspondence and accepted homography accuracy under the measured illumination and viewpoint changes?
\item How does global retrieval affect the candidate search needed for reliable place recognition, and when does its shortlist discard a geometrically useful reference?
\item How often does a plausible retrieval result lead to an accepted and correct place, and how does that differ from ground-truth planar alignment accuracy?
\item Does short contrastive adaptation of a compact projection head improve retrieval on held-out places and real cross-time images?
\item What accuracy, size and latency trade-offs follow from an executed descriptor export and INT8 quantisation attempt?
\end{enumerate}
Each question maps to an observable artifact. Pair statistics address the first; ranked predictions and candidate-count ablations address the second and third; the seeded training records address the fourth; native-versus-ONNX descriptors and timed inference address the fifth. Metric six-degree-of-freedom localisation is excluded from these answers because the downloaded evaluation subsets do not supply the required calibrated poses and three-dimensional map.

\section{Scope of the empirical contribution}
The public-data contribution comprises a deterministic HPatches subset and a custom Tokyo cross-time protocol. HPatches contributes real illumination and viewpoint changes with image-to-image homographies \citep{balntas2017}. Tokyo contributes photographs captured at different times with place and orientation labels \citep{torii2015}. The latter protocol uses a small reference gallery drawn from the query collection; it is deliberately named a cross-time subset rather than the official Tokyo street-view retrieval benchmark. Its results support a narrow comparison under matched labels, not a claim of benchmark-leading performance.

Synthetic queries serve a different purpose. They provide known geometric maps while varying brightness, gamma, colour balance, noise and contrast. Since all transformations are seeded, the same source identity can be followed through a controlled perturbation. However, altered pixels retain the original scene, visibility and much of its texture. Their performance is therefore interpreted as a stress test of an implementation, not as a substitute for real night imagery or evidence about winter conditions.

The methodological contribution is the integration and evaluation of available pretrained components in the Stitch setting. No new foundation model, feature detector, or neural matcher is claimed. The short training experiment changes only a projection head on frozen global descriptors. Its purpose is to expose the practical consequences of limited adaptation data and descriptor compression. Treating a neutral or negative outcome as a valid answer is necessary to avoid selecting a favourable training seed after observing the test set.

The software contribution includes modular extraction, matching, geometry, rendering, database and I/O modules; a command-line interface; configuration and tests; data/model acquisition scripts; evaluation and demonstration code; and a generated-report workflow. The reference database stores content alongside place identifiers so that the experiment exercises retrieval of an actual posted asset rather than merely drawing an arbitrary polygon after matching. Persisting that record on disk establishes local repeatability, while cross-device consistency remains future work.

\section{Assumptions and boundaries}
The reference photograph is assumed to represent a surface on which a planar overlay is meaningful. Within that model, an estimated homography maps reference pixels to query pixels. A street scene is generally not planar, so an accepted Tokyo homography is evidence of a consistent image region, not proof that every part of the street obeys the same map. This distinction motivates the separate reporting of geometric acceptance and true homography success.

The gallery is treated as known and locally available during the experiment. Reference descriptors and local features may be cached before a query arrives. Timings state whether they include this setup, image decoding, rendering and content-file access. The production design would need to constrain cache size and fetch latency. Those costs are analysed symbolically and through measured component sizes; they are not represented as measured cloud-service costs.

The work assumes that dataset access does not imply unrestricted redistribution or commercial permission. Software and imagery have separate provenance. Official documentation is retained where available, and manual-fetch instructions identify larger benchmarks whose image and mapping archives were not downloaded. A future product must resolve additional content rights, user consent and location privacy before reusing a research pipeline in public spaces.

\section{Organisation of the thesis}
\Cref{ch:background} reviews the representations, geometric models, datasets and anchor services relevant to the project. \Cref{ch:existing} analyses the original scripts and their recorded execution. \Cref{ch:architecture} describes the implemented modules and the proposed multi-user service boundary. \Cref{ch:datasets} defines the downloaded subsets and all preparation rules. \Cref{ch:methodology} develops the mathematical formulation and algorithms. \Cref{ch:setup} specifies hardware, metrics and experimental controls. \Cref{ch:results} reports the measured comparisons and validity threats. \Cref{ch:deployment} examines export, quantisation and scaling. \Cref{ch:ethics} treats privacy and abuse as system requirements. \Cref{ch:conclusion} answers the research questions and defines a staged roadmap. \ifFullReport The appendices provide reproduction commands, directory structure, qualitative evidence and key source listings.\else The appendix provides reproduction commands.\fi

\ifFullReport
The expanded technical treatment in \cref{ch:projective-foundations,ch:decision-analysis,ch:failure-analysis} derives coordinate transformations, separates decision metrics and traces stage-specific failures. \Cref{ch:hp-atlas,ch:tokyo-atlas} provide the complete scene-level evidence for every selected sequence and site. \Cref{ch:reliability,ch:persistent-transactions} document executable verification contracts and a proposed revision-aware service design. \Cref{app:pair-register} retains all oracle-pair outcomes, including failures. These additions expose more of the existing evidence without presenting it as a larger independently acquired experimental population.
\fi
